\section{The Theory in Action: Five Worked Examples}
\label{sec:examples}

The empirical demonstration in Section~\ref{sec:empirical} tests the theory's predictions against a historical corpus. But a theory of design form must also illuminate contemporary practice. This section applies the seven propositions to five concrete cases: one historical, two contemporary, and two emerging. The examples are \textit{illustrative}, not probative; they demonstrate the theory's explanatory reach and generate testable predictions, but they do not constitute independent empirical tests. Each example is described in factual detail and then read through the theoretical apparatus.

\subsection{The cantilever chair race, 1925--1928}

Between 1925 and 1928, four designers working in three cities independently converged on a shared morphological innovation: the cantilevered tubular steel chair. Marcel Breuer designed the Wassily (B3) in 1925 at the Bauhaus in Dessau, using cold-bent nickel-plated steel tubing inspired by his Adler bicycle. Mart Stam presented the first true cantilever (S33) at a Stuttgart housing exhibition in 1926, using straight gas-pipe segments joined at right angles. Ludwig Mies van der Rohe exhibited his cantilevered MR chair in 1927 at the Weissenhof exhibition, using a continuous curved tube that gave the seat its characteristic spring. Breuer returned with the Cesca (B32) in 1928, combining a cantilevered steel frame with a caned seat and back.

In the theory's terms, the material affordance of cold-bent steel tubing had been latent for decades (Proposition 2). Steel tubing was manufactured for bicycles and plumbing from the 1880s onward. The \textit{operation}; cold bending of thin-walled tubes to furniture scale; was the missing link. When Breuer recognized the bicycle handlebar as a furniture component, the forbidden region of the morphospace opened (Proposition 4). The convergence of four designers within three years, without direct collaboration on the initial innovations, confirms that the landscape attractor was real (Proposition 1): the fitness peak existed in the landscape before anyone climbed it.

The competitive dynamics among Breuer, Stam, and Mies constitute a poly-agentistic field (Proposition 6). Each designer navigated with a different cognitive cone: Breuer attended to visual lightness and Bauhaus pedagogy; Stam to structural minimalism and socialist housing; Mies to material elegance and spatial continuity. The resulting forms differ precisely along the axes each designer attended to. The lawsuit that followed (Stam vs.\ Breuer over priority) is itself evidence that independent agents converged on the same landscape feature from different directions.

\subsection{Neri Oxman's Silk Pavilion, MIT Media Lab, 2013}

The Silk Pavilion is a dome-shaped structure produced through a two-phase process. In the first phase, a CNC-controlled machine deposited a scaffold of silk threads across a steel frame, following a computationally generated pattern. In the second phase, 6,500 silkworms were placed on the scaffold and deposited additional silk over 20 days, guided by environmental parameters including light intensity and temperature gradients.

The result is a structure whose final form was not drawn by the designer. Neri Oxman and her team at the MIT Media Lab Mediated Matter group set the boundary conditions: scaffold geometry, species selection, environmental parameters. The biological agents navigated the form grammar at cellular scale, responding to local gradients of light and humidity. Each silkworm is a zero-order agent in the theory's sense (Proposition 5): it possesses a goal state (spinning a cocoon), a distance measure (proximity to suitable attachment points), and an adjustment mechanism (silk extrusion direction). The worm does not represent the dome; it responds to the local gradient.

This is multi-scale competence (Proposition 6) in its most literal form. The competence hierarchy runs from cellular silk production (millisecond timescale), through individual worm navigation (minute timescale), to collective pattern formation (day timescale), to the designer's scaffold specification (month timescale). Form arises in the intersection of these nested cognitive cones. The designer's care determines which axes are represented at the macro scale; the silkworms' biology handles the rest. The resulting form shows how far care reached at every scale: the scaffold's precision at the computational scale, the organic texture at the biological scale.

\subsection{Figma Weave and the designer as landscape navigator}

In 2025, Figma announced Weave, a generative AI feature embedded in its design tool. Weave produces interface layouts, component arrangements, and visual designs from textual or contextual prompts. The designer reviews, selects, modifies, and combines generated outputs rather than constructing each element manually.

In the theory's terms, the shape grammar $SG$ is now partially executed by a machine agent. The $C \to C$ partitions (generating new concept variations) are performed computationally at a speed and volume that exceed human manual capacity. The designer's role shifts from \textit{executing} grammar rules to \textit{steering the cognitive cone}: selecting which dimensions of the morphospace the system should attend to, which constraints to impose, which generated options to promote or discard.

This is a direct instantiation of Proposition 5. The designer's contribution is no longer primarily syntactic (performing form operations) but primarily attentional (setting the scope of care). Which axes does the designer actively represent? Accessibility, cultural context, emotional tone, brand consistency, interaction logic; these are the dimensions of the cognitive cone that the AI system does not autonomously navigate. The quality of the resulting design depends not on the AI's generative capacity but on the dimensionality of the human's care.

The theory predicts (Proposition 5) that a designer who narrows the cone to a single axis (e.g., visual novelty) will produce designs that fail along unrepresented axes (accessibility, usability, cultural sensitivity). The AI amplifies the designer's care; it does not replace it. If the cognitive cone is narrow, the AI generates variations within a narrow corridor. If the cone is wide, the AI explores a richer landscape. Weave makes the proposition observable in practice: the tool's output is a direct trace of the designer's active representational capacity.

\subsection{Chatbot interaction design}

When an interaction designer designs a conversational agent; a customer service chatbot, a medical triage assistant, a language tutoring system; she is designing a navigation algorithm for a non-human agent. The chatbot is an agent $A := (g, d, \delta)$ with a goal (task completion or user satisfaction), a distance measure (conversation state distance from resolution), and an adjustment mechanism (response generation policy).

The interaction designer is a \textbf{meta-agent} whose cognitive cone must encompass both the chatbot's cognitive cone and the user's cognitive cone (Proposition 6). The chatbot perceives the conversation through its training data and context window. The user perceives the interaction through expectations, emotional state, and communicative conventions. The designer must represent both sets of affordances simultaneously.

The quality of the interaction depends on how many axes the designer actively represents: task efficiency, user frustration, emotional tone, accessibility for screen readers, cultural appropriateness of phrasing, error recovery gracefully, edge case handling. Each unrepresented axis is an axis along which the design will fail (Proposition 5). A chatbot designed with care for task completion alone will be efficient but brittle; it will fail when the user is confused, upset, or operating in an unexpected context.

The theory's prediction is testable: designers who explicitly represent more axes during the design process (measurable through protocol analysis of design sessions) should produce chatbot interactions that score higher on multi-criteria evaluations. The form of the conversation; its rhythm, its recovery paths, its tone modulation; shows how far the designer's care reached.

\subsection{The Great Green Wall: continental-scale ecological design}

The Great Green Wall initiative, launched by the African Union in 2007, aims to restore 100 million hectares of degraded land across an 8,000-kilometer belt from Senegal to Djibouti. The project addresses desertification, food insecurity, climate change, and forced migration across the Sahel. As of 2024, approximately 18\% of the target area has been restored, with significant variation in approach and success across the eleven participating nations.

This is design at a scale that directly contradicts Le Corbusier's principle of man as the measure of all things. The ``form'' under design is not a single object but a multi-scalar configuration: soil microbiome composition (cellular-scale agent), root architecture and mycorrhizal networks (organism-scale agent), canopy structure and species diversity (community-scale agent), pastoral land-use patterns and community governance (social-scale agent), continental precipitation dynamics and carbon cycling (system-scale agent).

Each of these agents operates with its own cognitive cone on its own timescale, from microbial microseconds to climatic centuries. The theory predicts (Proposition 4) that collapsing the landscape to a single selection pressure; say, carbon sequestration; will produce configurations that fail under pressures the design was blind to. Early phases of the initiative focused heavily on tree planting (a single-pressure approach), and many plantations failed because they did not account for local soil conditions, pastoral mobility patterns, or indigenous species selection. The most successful restoration sites; such as the farmer-managed natural regeneration programs in Niger; employed a multi-pressure approach that coordinated biological, social, and economic agents simultaneously.

The theory frames this as a care problem (Proposition 7). The dimensionality of the cognitive cone must match the dimensionality of the pressure field. A design that represents only the carbon axis is an act of narrow care. A design that represents carbon, biodiversity, water table dynamics, pastoral rights, indigenous knowledge, and intergenerational equity is an act of expansive care. The form of the restored landscape shows how far care reached; and the failures show where it did not.
